% !TEX root = ../../../main.tex
% Appendix entry: preprocessing.surrogate -- "Surrogate generator" (Appendix A).
\entryfield{Algorithm}
\code{_run} in \code{Adapters/preprocessing_surrogate.py} creates one
\code{numpy.random.RandomState(seed)} and calls one of three functions in the same file on
the span $x_0,\dots,x_{n-1}$. This entry covers what the block does; how nulls are drawn
and judged across the app is in \cref{app:nulls}.

\begin{description}
\item[\texttt{phase\_randomize}] (\code{_phase_randomise}, a Fourier surrogate). Take
  $X = \mathrm{rfft}(x)$; draw $\phi_k \sim U[0, 2\pi)$ for every bin; set $\phi_0 = 0$ (the
  mean) and, for even $n$, the Nyquist bin's $\phi = 0$, so the inverse is real; return
  $\mathrm{irfft}(X_k e^{i\phi_k}, n)$. Every bin keeps its magnitude.
  \emph{Keeps}: the power spectrum exactly (so the autocorrelation) and the mean.
  \emph{Destroys}: every phase relation, so every local shape, any asymmetry in time, and the
  value distribution (the output tends towards Gaussian).
\item[\texttt{block\_shuffle}] (\code{_block_shuffle}). Refuse if \emph{block\_s} is 0 or
  $\textit{block\_s}\cdot f_s < 2$. With $L = \mathrm{round}(\textit{block\_s}\cdot f_s)$ and
  $B = \lfloor n/L \rfloor$ whole blocks: if $B \le 1$, return $x$ unchanged; otherwise
  reorder the $B$ blocks by a random permutation and leave the trailing partial block in
  place, so the length is exact.
  \emph{Keeps}: every sample value, and everything inside a block, so any motif shorter than a
  block. \emph{Destroys}: order and timing between blocks, and slow structure longer than a
  block; it adds a jump at each block seam.
\item[\texttt{aaft}] (\code{aaft_surrogate}; \citet{theiler1992testing}'s amplitude-adjusted
  Fourier transform). Rank-map $x$ onto sorted standard-normal draws; phase-randomise that
  Gaussian series with \code{_phase_randomise}; rank-map the result back onto $x$'s own
  sorted values. \emph{Keeps}: the exact multiset of values (the output is a permutation of
  $x$) and approximately the spectrum; the last step distorts it slightly (IAAFT, which
  iterates to fix that, is not offered). \emph{Destroys}: phase structure and shapes.
\end{description}

\emph{Seeds.} The seed is an ordinary parameter, so it is part of the step's recipe and of
its hash: the same recipe reproduces the same surrogate bit for bit, and a cached result is
reused only for the same seed. When a run's null is drawn (\code{draw_paired_null} and
\code{run_paired_recipe} in \code{Working/run_groups.py}), the chain is copied with this
block prepended, and draw $i$ runs at seed $\text{seed} + i$ (base seed 0 unless set), each
as its own run linked to the real one by \code{runs.surrogate_of_run_id}. For block
shuffling, unless the caller fixes \emph{block\_s}, the null sets it to twice the longest
span the real run detected (at least two samples), and skips the null with a reason when the
real run detected nothing (\code{surrogate_draw_params}).

\entryfield{Assumptions}
\begin{itemize}
\item \emph{Units.} None needed: every method returns values in the input's unit (AAFT
  returns the input's own values).
\item \emph{Sampling rate.} Only \texttt{block\_shuffle} uses $f_s$, to turn \emph{block\_s}
  seconds into samples. The output keeps the input's $f_s$.
\item \emph{Missing values.} Not handled. \texttt{phase\_randomize}: one NaN makes the FFT,
  and so the whole output, NaN. \texttt{block\_shuffle}: NaNs move with their blocks.
  \texttt{aaft}: NaNs sort last, so they are reassigned to the positions where the
  randomised series is highest, scattered through the output.
\item \emph{Stationarity.} Phase randomisation and AAFT treat the span as one stationary
  stretch: a slow trend becomes part of the spectrum and is smeared, and the FFT treats the
  span as periodic, so a level difference between its two ends leaks into every frequency.
\end{itemize}

\entryfield{Outputs}
A \texttt{Signal} of the same length and $f_s$. No \code{meta}: the method and seed are in
the step's recipe.

\entryfield{Cost}
\texttt{phase\_randomize}: one forward and one inverse real FFT, $O(n \log n)$. \texttt{aaft}:
two sorts and the same FFTs, $O(n \log n)$. \texttt{block\_shuffle}: $O(n)$. Memory a few
copies of the span. Each null draw reruns the whole chain after it, so the chain's cost, not
this block's, dominates.

\entryfield{Failure modes}
\begin{itemize}
\item \emph{A silent identity null.} Block shuffling on a span shorter than two blocks
  returns the input unchanged with no warning; a chain then ``beats'' nothing and matches
  its null exactly.
\item \emph{NaN} in the span (above): a phase-randomised null of a span with a gap is all
  NaN.
\item \emph{The wrong null.} A phase-randomised copy keeps the spectrum: a smooth slow
  shape lives in the spectrum, so a detector of smooth slow shapes can find them in the
  copies too. The blocks that need a different null (the quiet-stretch null of
  \blockref{interrogation.event_families} and \blockref{interrogation.cycle_shapes}) build
  their own.
\item \emph{A non-Gaussian measure against a Gaussian null.} Phase randomisation makes the
  values roughly Gaussian; a measure that reads the value distribution (multiscale entropy
  with a MAD-based tolerance) beat its phase-randomised surrogates 20 of 20 times on a
  non-Gaussian linear process (module docstring). Use \texttt{aaft} for such a measure.
\item \emph{Different generators.} The block uses \code{RandomState}; the older pipeline
  in \code{Pipelines/drop_motifs/nulls1.py} uses \code{default_rng}. The same seed gives
  different surrogates in the two places.
\end{itemize}

\entryfield{Validation}
\code{tests/test_surrogate.py} checks that a fixed seed reproduces a bit-identical surrogate
for every method, that phase randomisation keeps every rFFT magnitude (to $10^{-6}$), that
AAFT's output is a permutation of its input with its spectrum within a factor of 2 per band,
and that the AAFT moved here is bit-identical to the old one.
\code{tests/test_null_draws_methods_cut.py} checks the refusal of a block under two samples
or of no block length, and that a paired run draws $N$ surrogates at $N$ seeds.
\code{tests/test_drop_motifs_nulls1.py} checks, through \code{nulls1}'s wrapper, that block
shuffling keeps the length and every sample, actually reorders, and keeps a block's
waveform intact.

\entryfield{References}
\citet{theiler1992testing} (phase randomisation and AAFT). Nulls in general: \cref{app:nulls}.

\entryfield{Tests}
\code{tests/test_surrogate.py}, \code{tests/test_null_draws_methods_cut.py},
\code{tests/test_drop_motifs_nulls1.py}, \code{tests/test_discovery_seeded_search.py} (the
methods the seed search offers are the block's), \code{tests/test_band_scope.py},
\code{tests/test_hpc_job_dependencies.py}, \code{tests/test_adapter_spec.py}.
